\documentclass[11pt]{article}       % set main text size
\usepackage[letterpaper,            % set paper size to letterpaper. change to a4paper for resumes outside of North America
top=0.5in,                          % specify top page margin
bottom=0.5in,                       % specify bottom page margin
left=0.5in,                         % specify left page margin
right=0.5in]{geometry}              % specify right page margin
                       
\usepackage{XCharter}               % set font
\usepackage[T1]{fontenc}            % output encoding
\usepackage[utf8]{inputenc}         % input encoding
\usepackage{enumitem}               % enable lists for bullet points: itemize and \item
\usepackage[hidelinks]{hyperref}    % format hyperlinks
\usepackage{titlesec}               % enable section title customization
\raggedright                        % disable text justification
\pagestyle{empty}                   % disable page numbering

% ensure PDF output will be all-Unicode and machine-readable
\input{glyphtounicode}
\pdfgentounicode=1

% format section headings: bolding, size, white space above and below
\titleformat{\section}{\bfseries\large}{}{0pt}{}[\vspace{1pt}\titlerule\vspace{-6.5pt}]

% format bullet points: size, white space above and below, white space between bullets
\renewcommand\labelitemi{$\vcenter{\hbox{\small$\bullet$}}$}
\setlist[itemize]{itemsep=-2pt, leftmargin=12pt, topsep=6pt} %%% Test various topsep values to fix vertical spacing errors

% resume starts here
\begin{document}

% name
\centerline{\Huge Nick Graham}

\vspace{5pt}

% contact information
\centerline{\href{mailto:nick@nickjgraham.com}{nick@nickjgraham.com} | \href{https://nickjgraham.com}{nickjgraham.com} | \href{https://linkedin.com/in/ngraham2/}{linkedin.com/in/ngraham2/}}

\vspace{-10pt}

% skills section
\section*{Skills}

\textbf{Application Servers:} Apache HTTPd, IIS, Nginx, Ping Federate (SAML/Oauth), Tomcat \\

\textbf{Automation:} Ansible, Artifact Repositories, Jenkins, Packer, Terraform \\

\textbf{Containerization:} Docker, Kubernetes, Rancher \\

\textbf{Hypervisors:} Amazon Web Services, VMware \\

\textbf{Operating Systems:} Red Hat Linux, Ubuntu, Windows Server \\

\textbf{Programming:} Bash, Javascript, Powershell, Python \\


\vspace{-6.5pt}

% experience section
\section*{Experience}

\textbf{Principal Infrastructure Engineer,} {Capital District Physician's Health Plan} -- Albany, NY \hfill 2016 - Present \\
\vspace{-9pt}
\begin{itemize}
  
  \item Architected and executed Single Sign-On (SSO) system expansion to AWS, creating a multi-region, highly available system with no downtime since its creation. This greatly improved the reliability of our authentication system for both employees and customers.
  
  \item Fully automated deployment and configuration management of Mulesoft application server architecture using Ansible, including over 80 individual services across 4 SDLC environments. This allowed for consistent, reliable, quick deployments of new services and updates.
  
  \item Led the containerization of legacy Tomcat services into fully automated, reliable solution on Kubernetes. This enabled developers to deploy new services with minimal effort, and provided a consistent, reliable platform for those services to run on.
  
  \item Created many CI/CD pipelines in Jenkins to turn complex processes like code builds, file deployments, or configuration changes into scheduled, or one-click operations. This greatly reduced the time and effort required to deploy changes, as well as reducing the chance for human error.
  
  \item Automated the transition from CentOS to RHEL, and subsequent upgrades to RHEL 9. This reduced the amount of labor required to perform these upgrades by the whole engineering staff.
  
  \item Organized and streamlined Ansible playbook usage, including the development of easy-to-use environment build scripts, and personal training/mentoring for other engineers. This enhanced the use of Ansible and improved automation capabilities for everyone's work across the team.
  
  \item Created a custom internal site with Python that integrates with our SSO system APIs. This provides users across the company with a central place to find SSO links, rather than relying on links maintained in wikis or emails.
  
\end{itemize}

\textbf{Systems Engineer,} {Xerox Corporation} -- Albany, NY \hfill 2013 - 2016 \\
\vspace{-9pt}
\begin{itemize}
  
  \item Created a VM deployment process automation from the ground up with VMware Orchestrator.
  
  \item Architected and implemented up all of the backup infrastructure required to support disaster recovery for all of Xerox's systems
  
  \item Enhanced the system lifecycle process around all of our VMware infrastructure. Everything from creation of new servers, patching existing ones, updating configurations, and server decommissioning was streamlined and automated in some way.
  
  \item Led VMware expansion to two new datacenters in China, including hardware purchasing, deployment, and configuration.
  
\end{itemize}

\textbf{Information Systems Intern,} {Saratoga Hospital} -- Saratoga Springs, NY \hfill 2013 - 2013 \\
\vspace{-9pt}
\begin{itemize}
  
  \item Created custom code for the ticket tracking system (Spiceworks), in order to add some custom fields and workflows that the base product didn't provide.
  
  \item Created a fully featured PC deployment solution using the Microsoft Deployment Toolkit (MDT). Before my arrival PC deployments were all done manually, clicking through the installer. After deploying MDT new PC builds were as simple as starting a machine and selecting the PXE boot option.
  
\end{itemize}

\textbf{IT Technician,} {SUNY Canton} -- Canton, NY \hfill 2009 - 2013 \\
\vspace{-9pt}
\begin{itemize}
  
  \item Created a PHP application to track the association between switches, patch panel ports, and wall ports in classrooms. Identifying where various switch ports went was a huge hassle at the time, and my contribution removed the painful process of having to manually trace out each network connection when making port changes.
  
\end{itemize}


\vspace{-18.5pt}

% projects section
\section*{Projects}

\textbf{Resume Website} \hfill \href{https://nickjgraham.com}{nickjgraham.com} \\
\vspace{-9pt}
\begin{itemize}
  
  \item I created a website to host my resume in a more unique way than a simple PDF.
  
  \item It's written in Python, runs on Kubernetes, and includes an about page with more information on how it works.
  
\end{itemize}

\textbf{Home Lab} \hfill  \\
\vspace{-9pt}
\begin{itemize}
  
  \item I run a small server lab in my home, which allows me to test new technologies and stay current with skills that I may not use in my day-to-day work.
  
  \item The lab consists of one physical server with around 60 virtual machines running a variety of Windows and Linux based services.
  
\end{itemize}


\vspace{-18.5pt}

% education section
\section*{Education}

\textbf{SUNY Canton} -- BS, Information Technology \hfill 2013 \\


\end{document}